\section{The original Rees map on its actual original ratio tuple}
\begin{definition}[The same original ratio tuple]
\label{def:whole-original-fraction-ratio-tuple}
\lean{EnvelopeProperty.BlowupHomogeneousTautologicalLine.originalFractionRatioTuple}
\leanok
Let $A$ be a domain, $I\subset A$ an ideal, and
$g=(g_0,\ldots,g_m)$ its specified original generating family.
On the actual native fraction chart $S=A[I/g_i]$, define
$a_j=g_j/g_i\in S$ using the original ideal membership of $g_j$.
This tuple has exactly the original $m+1$ coordinates.
\end{definition}
\begin{proof}
The specified ideal equality implies $g_j\in I$ for every original index $j$.
Each displayed quotient is therefore one of the defining generators of the
native subalgebra $S\subset\operatorname{Frac}(A)$, and hence is an element
of $S$. These elements define the stated function on the original index set.
The native definition uses the fraction field's total division; the actual
geometric chart inclusion below is used only under $g_i\ne0$.
\end{proof}
\begin{lemma}[The original pivot is one]
\label{lem:whole-original-fraction-ratio-pivot}
\uses{def:whole-original-fraction-ratio-tuple}
\lean{EnvelopeProperty.BlowupHomogeneousTautologicalLine.originalFractionRatioTuple_pivot}
\leanok
If the specified original $g_i$ is nonzero, then $a_i=1$ in $S$.
\end{lemma}
\begin{proof}
The injection $A\hookrightarrow\operatorname{Frac}(A)$ carries $g_i$ to a
nonzero element. Therefore its quotient by itself is one in
$\operatorname{Frac}(A)$. The defining inclusion $S\hookrightarrow\operatorname{Frac}(A)$
is injective, giving the asserted identity in the actual chart ring.
\end{proof}
\begin{lemma}[The actual ratio tuple spans the unit ideal]
\label{lem:whole-original-fraction-ratio-span}
\uses{lem:whole-original-fraction-ratio-pivot}
\lean{EnvelopeProperty.BlowupHomogeneousTautologicalLine.originalFractionRatioTuple_span}
\leanok
Under the same nonzero-pivot hypothesis, $(a_0,\ldots,a_m)=S$.
\end{lemma}
\begin{proof}
The ideal generated by the tuple contains its $i$th entry. By
\ref{lem:whole-original-fraction-ratio-pivot}, that entry is one, so the
ideal is the unit ideal. No unit-span certificate is an input.
\end{proof}
\begin{proposition}[The whole original Rees restriction is the same tuple map]
\label{prop:whole-original-rees-fraction-ratio-map}
\uses{lem:whole-original-fraction-ratio-span,thm:resume-blowup-affine-generator-chart-readback,thm:fc-generated-tuple-unit-chart}
\lean{EnvelopeProperty.BlowupHomogeneousTautologicalLine.affineGeneratorMap_eq_originalFractionRatioTuple}
\leanok
Let $r:R\to A$ be the original coefficient map. On the actual original native
$i$th fraction chart, the specified-generator original Rees morphism
$q:\operatorname{Bl}_I(\operatorname{Spec}A)\to\mathbf P_R^m$ satisfies
$c_iq=\operatorname{Tuple}((A\to S)r,(a_0,\ldots,a_m))$.
Both sides have the same original target $\mathbf P_R^m$.
\end{proposition}
\begin{proof}
The proved original Rees chart readback identifies $c_iq$ with
polynomial evaluation on the original target's $i$th coordinate chart: the
coefficient map is exactly $(A\to S)r$, and each original dehomogenized
coordinate $X_j/X_i$ maps to the actual ratio $g_j/g_i=a_j$.
By \ref{lem:whole-original-fraction-ratio-span}, the actual ratio tuple defines
its native tuple morphism, and by \ref{lem:whole-original-fraction-ratio-pivot}
its original $i$th coordinate equals one. The native tuple chart evaluation
theorem then gives precisely the same polynomial evaluation followed by the
same original coordinate-chart inclusion. The two entire scheme morphisms
are equal by transitivity. Neither map equality, section matching, a whole
line comparison, nor a padded replacement target has been assumed.
\end{proof}
